\documentclass{fiche}
\metadonnees{7}{Apprendre en secret}{Bloc 2 — Ce qui est possible : futur et modaux}

\begin{document}
\entete

\objectifs{%
\textbf{Langue} : \emph{can}, \emph{could}, \emph{be able to} ; permission et interdiction ;
lexique de l'apprentissage.\par
\textbf{Texte} : Frederick Douglass, \emph{Narrative of the Life of Frederick Douglass}
(1845), chapitre \textsc{vii}.\par
\textbf{Durée} : 1\,h\,30 — exercices 35 min, texte et questions 30 min, version 25 min.}

%% ===================================================================
\exercice[13 min]{\emph{Can}, \emph{could}, \emph{be able to}}

\rappel{\textbf{Rappel.} \emph{Can} dit la capacité au présent, \emph{could} au passé ou
dans l'hypothèse. \emph{Be able to} le remplace là où le modal n'a pas de forme : infinitif,
futur, present perfect.}

\consigne{Complétez avec la forme qui convient. Plusieurs réponses sont parfois possibles.}

\begin{anglais}
\begin{enumerate}[label=\arabic*.,itemsep=2.2mm,leftmargin=*]
  \item At seven he \trou{could} already spell words of three or four letters.
  \item He would like to \trou{be able to} read the newspapers.
  \item The house was burning, but he \trou{was able to} get out through the window.
    \emph{(réussite ponctuelle)}
  \item I \trou{have never been able to} understand that argument.
  \item ``\trou{Can} you read?'' he asked the boy.
  \item By next year the child \trou{will be able to} read as well as his masters.
  \item She \trou{could} hear the men singing on the quays every morning.
  \item Nobody \trou{could} have foreseen what the boy would become.
  \item After three months of secret lessons he \trou{managed to} read a whole page aloud.
  \item If I \trou{could} choose, I would go to sea.
  \item He \trou{was not able to} finish the book that night. \emph{(ou} could not\emph{)}
  \item They dared not speak, but they \trou{could} look at each other.
\end{enumerate}
\end{anglais}

\note{Phrases 3 et 9 : pour une réussite unique et ponctuelle dans le passé, l'anglais dit
\emph{was able to} ou \emph{managed to}, jamais \emph{could}. \emph{Could} ne vaut au passé
que pour une capacité durable (1, 7, 12). À la forme négative, en revanche, \emph{could not}
convient dans les deux cas (11).}

\note[Formes manquantes]{Un modal n'a ni infinitif, ni participe, ni futur : on passe alors
par \emph{be able to} (2, 4, 6). \emph{Could have} + participe (8) sert à l'irréel du passé :
\og personne n'aurait pu \fg.}

%% ===================================================================
\exercice[10 min]{Permission et interdiction}

\rappel{\textbf{Rappel.} \emph{May} et \emph{can} donnent la permission, \emph{be allowed to}
la reformule à tous les temps. \emph{Must not} interdit, \emph{need not} dispense.}

\consigne{Complétez avec \emph{may}, \emph{can}, \emph{be allowed to}, \emph{must not} ou
\emph{need not}.}

\begin{anglais}
\begin{enumerate}[label=\arabic*.,itemsep=2.2mm,leftmargin=*]
  \item In that state, a slave \trou{was not allowed to} learn to read.
  \item ``You \trou{must not} tell anyone where you got the book.''
  \item \trou{May} I take the newspaper with me on my errand?
  \item You \trou{need not} bring bread today; there is enough in the house.
  \item The boys \trou{were allowed to} play in the street until dusk.
  \item Visitors \trou{must not} walk on the lawn.
  \item ``You \trou{can} borrow it, but you \trou{must not} lose it.''
  \item Since last year we \trou{have been allowed to} use the library.
  \item You \trou{need not} come if you are tired.
  \item He was forbidden to read; he dared not ask whether he \trou{might} keep the book.
\end{enumerate}
\end{anglais}

\note[Deux distinctions]{\emph{Must not} n'est pas la négation de \emph{must} : c'est une
interdiction (\og il ne faut pas \fg). L'absence d'obligation se dit \emph{need not} ou
\emph{do not have to} (\og ce n'est pas la peine \fg). Phrase 10 : au discours indirect
passé, \emph{may} devient \emph{might}.}

%% ===================================================================
\exercice[12 min]{Apprendre, interdire}
\consigne{a) Associez chaque mot à sa traduction.}

\begin{multicols}{2}
\begin{enumerate}[label=\arabic*.,itemsep=1.2mm,leftmargin=*]
  \item to forbid \hfill \trou{d}
  \item unlawful \hfill \trou{g}
  \item an errand \hfill \trou{b}
  \item to bestow \hfill \trou{k}
  \item an urchin \hfill \trou{a}
  \item to spell \hfill \trou{i}
  \item to utter \hfill \trou{f}
  \item a purpose \hfill \trou{l}
  \item to dread \hfill \trou{c}
  \item to shun \hfill \trou{j}
  \item wretched \hfill \trou{e}
  \item a curse \hfill \trou{h}
\end{enumerate}
\columnbreak
\begin{enumerate}[label=\alph*.,itemsep=1.2mm,leftmargin=*]
  \item un galopin, un gamin des rues
  \item une course, une commission
  \item redouter
  \item interdire
  \item misérable, pitoyable
  \item proférer, prononcer
  \item illégal
  \item une malédiction
  \item épeler
  \item fuir, éviter
  \item octroyer, faire don de
  \item un but, un dessein
\end{enumerate}
\end{multicols}

\note[Deux préfixes]{\emph{Un-} nie l'adjectif (\emph{unlawful}, \emph{unmanageable},
\emph{unpardonable}, \emph{unabated}) ; c'est le préfixe privatif le plus productif de
l'anglais. \emph{To forbid} suit le schéma \emph{bid, bade, bidden} : \emph{forbid, forbade,
forbidden}.}

\consigne{b) Complétez avec un mot de la liste.}
\begin{anglais}
\begin{enumerate}[label=\arabic*.,itemsep=2mm,leftmargin=*]
  \item His master \trou{forbade} his wife to teach him, because it was \trou{unlawful}.
  \item When he was sent on an \trou{errand}, he always took his book with him.
  \item He gave his bread to the hungry little \trou{urchins} of the street.
  \item What his master most \trou{dreaded}, he himself most desired.
  \item He set out with high hope and a fixed \trou{purpose}.
\end{enumerate}
\end{anglais}

%% ===================================================================
\newpage
\section{Texte}

\noindent\small\textit{Né esclave dans le Maryland, Frederick Douglass est envoyé enfant à
Baltimore. Sa maîtresse a commencé à lui apprendre l'alphabet ; son maître le lui a interdit,
au motif qu'instruire un esclave le rendrait \og ingouvernable \fg{} et le rendrait à jamais
impropre à l'esclavage. Douglass en tire aussitôt la conclusion inverse.}\normalsize

\begin{texte}
The plan which I adopted, and the one by which I was most successful, was that of making
friends of all the little white boys whom I met in the street. As many of these as I could, I
converted into teachers. With their kindly aid, obtained at different times and in different
places, I finally succeeded in learning to read. When I was sent of
\gl[an errand]{errands}{une course, une commission}, I always took my book with me, and by
going one part of my errand quickly, I found time to get a lesson before my return. I used
also to carry bread with me, enough of which was always in the house, and to which I was
always welcome; for I was much better off in this regard than many of the poor white children
in our neighborhood. This bread I used to \gl[to bestow]{bestow}{octroyer, faire don de} upon
the hungry little \gl[an urchin]{urchins}{un galopin, un gamin des rues}, who, in return,
would give me that more valuable bread of knowledge. I am strongly tempted to give the names
of two or three of those little boys, as a testimonial of the gratitude and affection I bear
them; but prudence forbids --- not that it would injure me, but it might embarrass them; for
it is almost an \gl[unpardonable]{unpardonable}{impardonnable} offence to teach slaves to
read in this Christian country. It is enough to say of the dear little fellows, that they
lived on Philpot Street, very near Durgin and Bailey's \gl[a ship-yard]{ship-yard}{un
chantier naval}. I used to talk this matter of slavery over with them. I would sometimes say
to them, I wished I could be as free as they would be when they got to be men. ``You will be
free as soon as you are twenty-one, \emph{but I am a slave for life!} Have not I as good a
right to be free as you have?'' These words used to trouble them; they would express for me
the liveliest sympathy, and console me with the hope that something would occur by which I
might be free.

I was now about twelve years old, and the thought of being \emph{a slave for life} began to
bear heavily upon my heart. Just about this time, I got hold of a book entitled ``The
Columbian \gl[an orator]{Orator}{un orateur}.'' Every opportunity I got, I used to read this
book. Among much of other interesting matter, I found in it a dialogue between a master and
his slave. The slave was represented as having run away from his master three times. The
dialogue represented the conversation which took place between them, when the slave was
\gl[to retake]{retaken}{reprendre, rattraper} the third time. In this dialogue, the whole
argument in behalf of slavery was brought forward by the master, all of which was
\gl[to dispose of]{disposed of}{réfuter, régler le sort de} by the slave. The slave was made
to say some very \gl[smart]{smart}{fin, habile} as well as impressive things in reply to his
master --- things which had the desired though unexpected effect; for the conversation
resulted in the voluntary \gl[emancipation]{emancipation}{l'affranchissement} of the slave on
the part of the master.

In the same book, I met with one of Sheridan's mighty speeches on and in behalf of Catholic
emancipation. These were choice documents to me. I read them over and over again with
\gl[unabated]{unabated}{sans faiblir} interest. They gave tongue to interesting thoughts of
my own soul, which had frequently flashed through my mind, and died away for want of
\gl[utterance]{utterance}{l'expression, la parole}. The moral which I gained from the dialogue
was the power of truth over the conscience of even a slaveholder. What I got from Sheridan
was a bold \gl[a denunciation]{denunciation}{une dénonciation} of slavery, and a powerful
\gl[a vindication]{vindication}{une défense, une revendication} of human rights.
\end{texte}
\source{Frederick Douglass, \emph{Narrative of the Life of Frederick Douglass, an American
Slave}, Boston, Anti-Slavery Office, 1845, ch.~\textsc{vii}.}

\partiel[15 min]{Questions}
\consigne{\en{Answer in English, in full sentences.}}

\begin{enumerate}[label=\arabic*.,itemsep=2mm,leftmargin=*]
  \item \en{What plan did Douglass adopt, and why did it work?}
    \lignes{2}
    \corr{He made friends of the little white boys he met in the street and turned as many of
    them as he could into teachers. It worked because the lessons were short, scattered and
    unnoticed.}
  \item \en{Explain the exchange of bread. What are the two kinds of bread?}
    \lignes{3}
    \corr{He carried bread from his master's house, where there was always enough, and gave
    it to hungry white children who were poorer than he was; in return they gave him
    \emph{that more valuable bread of knowledge}, that is, reading lessons.}
  \item \en{Why does he refuse to give the boys' names?}
    \lignes{2}
    \corr{Not to protect himself, but to protect them: teaching a slave to read is an almost
    unpardonable offence, and naming them might get them into trouble.}
  \item \en{What argument does he put to the boys, and how do they react?}
    \lignes{3}
    \corr{He points out that they will be free at twenty-one while he is a slave for life,
    and asks whether he has not as good a right to freedom as they have. The words trouble
    them; they pity him and hope something will free him.}
  \item \en{What does he find in \emph{The Columbian Orator}, and what does he learn from
    it?}
    \lignes{3}
    \corr{A dialogue in which a recaptured slave answers every argument of his master so well
    that the master frees him, and a speech by Sheridan for Catholic emancipation. He learns
    the power of truth over a slaveholder's conscience, and finds words for thoughts he
    already had.}
  \item \en{``They gave tongue to interesting thoughts of my own soul.'' Why does reading
    matter so much to him?}
    \lignes{3}
    \corr{Because the thoughts were already his but had no words and \emph{died away for want
    of utterance}. Reading does not give him ideas; it gives him the means to say them, and
    therefore to argue.}
\end{enumerate}

\note[Les modaux dans le texte]{\emph{As many of these as I could}, \emph{I wished I could be
as free}, \emph{something would occur by which I might be free} : capacité, souhait,
possibilité. Noter \emph{I used to carry}, \emph{they would express} : l'habitude passée de
la fiche 2, ici sur toute une page.}

%% ===================================================================
\newpage
\section{Version}
\consigne{Traduisez le passage en français. Il suit immédiatement le texte.}

\begin{version}
The more I read, the more I was led to \gl[to abhor]{abhor}{abhorrer, exécrer} and detest my
enslavers. I could regard them in no other light than a band of successful robbers, who had
left their homes, and gone to Africa, and stolen us from our homes, and in a strange land
reduced us to slavery. I loathed them as being the meanest as well as the most wicked of men.
As I read and contemplated the subject, behold! that very discontentment which Master Hugh
had predicted would follow my learning to read had already come, to torment and sting my soul
to unutterable \gl[anguish]{anguish}{l'angoisse, le tourment}. As I
\gl[to writhe]{writhed}{se tordre (de douleur)} under it, I would at times feel that learning
to read had been a curse rather than a blessing. It had given me a view of my
\gl[wretched]{wretched}{misérable} condition, without the remedy. It opened my eyes to the
horrible pit, but to no ladder upon which to get out. In moments of agony, I envied my
fellow-slaves for their stupidity. I have often wished myself a beast. I preferred the
condition of the meanest reptile to my own. Any thing, no matter what, to get rid of
thinking! It was this everlasting thinking of my condition that tormented me. There was no
getting rid of it. It was pressed upon me by every object within sight or hearing, animate or
inanimate. The silver \gl[a trump]{trump}{une trompette (sens ancien et biblique)} of freedom
had roused my soul to eternal wakefulness. Freedom now appeared, to disappear no more
forever. It was heard in every sound, and seen in every thing. It was ever present to torment
me with a sense of my wretched condition. I saw nothing without seeing it, I heard nothing
without hearing it, and felt nothing without feeling it. It looked from every star, it smiled
in every calm, breathed in every wind, and moved in every storm.
\end{version}
\source{\emph{Ibid.}}

\lignes{22}

\corr{\textbf{Proposition de traduction.} Plus je lisais, plus j'en venais à abhorrer et à
détester ceux qui m'asservissaient. Je ne pouvais les voir autrement que comme une bande de
voleurs heureux, qui avaient quitté leur pays, gagné l'Afrique, nous avaient arrachés à nos
foyers et, sur une terre étrangère, réduits en esclavage. Je les haïssais comme les plus vils
autant que les plus méchants des hommes. À mesure que je lisais et que je méditais là-dessus,
voici que ce mécontentement même que maître Hugh avait annoncé si j'apprenais à lire était
déjà venu tourmenter mon âme et la piquer d'une angoisse indicible. Tandis que je me tordais
sous elle, je sentais par moments que savoir lire avait été une malédiction plutôt qu'une
bénédiction. Cela m'avait donné la vue de ma condition misérable, sans le remède. Cela
m'avait ouvert les yeux sur la fosse horrible, mais sans échelle pour en sortir. Dans mes
moments d'angoisse, j'enviais à mes compagnons d'esclavage leur stupidité. J'ai souvent
souhaité être une bête. Je préférais la condition du plus vil reptile à la mienne. N'importe
quoi, n'importe quoi, pourvu que je cesse de penser ! C'était cette pensée éternelle de ma
condition qui me tourmentait. Il n'y avait pas moyen de m'en défaire. Tout objet à portée de
vue ou d'oreille me l'imposait, animé ou inanimé. La trompette d'argent de la liberté avait
éveillé mon âme à une veille éternelle. La liberté venait d'apparaître, pour ne plus jamais
disparaître. On l'entendait dans chaque son, on la voyait en toute chose. Elle était sans
cesse là pour me tourmenter du sentiment de ma misère. Je ne voyais rien sans la voir,
je n'entendais rien sans l'entendre, je ne sentais rien sans la sentir. Elle regardait de
chaque étoile, souriait dans chaque accalmie, respirait dans chaque vent et se mouvait dans
chaque tempête.}

\note[Choix de traduction 1]{\emph{The more I read, the more I was led to\ldots} : la double
comparaison progressive se rend mot pour mot (\og plus\ldots plus\ldots \fg), sans article et
sans \emph{que}. Attention à l'ordre : en anglais le comparatif est en tête de chaque membre,
en français aussi.}

\note[Choix de traduction 2]{\emph{I could regard them in no other light than\ldots} :
\emph{in this light} = \og sous ce jour \fg. La négation porte sur \emph{no other} : \og je
ne pouvais les voir autrement que comme \fg.}

\note[Choix de traduction 3]{\emph{that very discontentment which Master Hugh had predicted
would follow my learning to read} : deux relatives emboîtées et un verbe rejeté très loin.
Reconstruire : \emph{Master Hugh had predicted [that discontentment would follow my learning
to read]}. \emph{My learning to read} est un gérondif avec sujet possessif : \og le fait que
j'apprenne à lire \fg.}

\note[Choix de traduction 4]{\emph{but to no ladder upon which to get out} : l'ellipse est
violente ; le français doit rétablir (\og mais sans échelle pour en sortir \fg). La
préposition antéposée à la relative (\emph{upon which}) appartient à la langue écrite : à
l'oral on dirait \emph{no ladder to get out on}.}

\note[Choix de traduction 5]{\emph{I saw nothing without seeing it} : \emph{without} + gérondif
se rend par \og sans \fg{} + infinitif. La triple reprise est un effet oratoire hérité des
sermons : la conserver intacte.}

%% ===================================================================
\partiel{Lexique à mémoriser}
\begin{itemize}[label=--,itemsep=0.8mm,leftmargin=*]\small
  \item \textbf{to forbid} — interdire ; forbade, forbidden
  \item \textbf{unlawful} — illégal ; un- + law + -ful
  \item \textbf{an errand} — une course, une commission
  \item \textbf{to succeed in} — réussir à ; \textit{noms} success, \textit{et} succession \textit{« la suite »}
  \item \textbf{to bestow} — octroyer
  \item \textbf{a fellow} — un garçon, un type ; fellow-slaves \textit{« compagnons d'esclavage »}
  \item \textbf{to bear} — porter ; \textit{ici} to bear somebody affection
  \item \textbf{to injure} — nuire à, blesser ; \textit{nom} an injury
  \item \textbf{an offence} — une infraction ; \textit{verbe} to offend
  \item \textbf{in behalf of} — en faveur de ; on behalf of \textit{« au nom de »}
  \item \textbf{to bring forward} — avancer, présenter (un argument)
  \item \textbf{to dispose of} — se débarrasser de, réfuter
  \item \textbf{smart} — fin, habile ; \textit{et non « élégant » au sens français}
  \item \textbf{to utter} — proférer ; \textit{noms} utterance\textit{, adjectif} utter \textit{« complet »}
  \item \textbf{to loathe} — exécrer \textit{(fiche 3)}
  \item \textbf{mean} — vil, mesquin ; \textit{superlatif} the meanest
  \item \textbf{wicked} — méchant ; \textit{se prononce en deux syllabes}
  \item \textbf{a curse} — une malédiction ; \textit{contraire} a blessing
  \item \textbf{wretched} — misérable ; \textit{nom} a wretch \textit{« un malheureux »}
  \item \textbf{to get rid of} — se débarrasser de
\end{itemize}

\end{document}
